\section{Triggers and the Alert Bank}
\label{sec:triggers}

Who decides when the model is consulted, and where does the text come from? Both
are infrastructure rather than contribution, and we say so, because
event-triggered invocation is already well covered~\cite{event-triggered-llm}.
What matters is that the infrastructure be shared exactly enough that no
comparison can be confounded by it.

A proposal opportunity opens when an operational alert arrives or when a
calibrated telemetry detector crosses its threshold. Six rules exist: alert only,
CUSUM~\cite{page-cusum}, Page-Hinkley, their disjunction with the alert channel,
which is the primary rule, and two budget-matched schedules. Detectors are
stateless about suppression: one emits a raw candidate and two wrappers decide,
composed as \texttt{MaxProposals(Refractory(}\textit{d}\texttt{, 5), 2)}, so the
window and the cap are uniform across rules, including rules written after the
design was frozen. Thresholds come from Monte Carlo on development and
calibration nulls only, at the nearest-rank 95th percentile of the per-episode
maximum statistic over 60 stationary twin paths, so three of 60 null episodes
fire by construction. The window of 5 periods covers the longest registered pulse
and pause, and the cap of 2 is the value $\alpha_j = \alpha\,2^{-j}$ was designed
around. Two properties matter for the claims. The demand-side detectors read only the
previous demand and the alert channel, both exogenous to the policy, so a trigger
trace is a pure function of the episode's demand and alert streams and two arms
sharing a rule must produce bitwise-identical traces, which the suite asserts.
Without that, a comparison between two arms would conflate a different invocation
pattern with a different policy. And a denylist, enforced at runtime on
every feature mapping and statically by a syntax-tree scan, forbids future demand
and arrivals, pattern names, instance identifiers, split and test labels and every
field of the hidden incident record. It is deliberately over-broad on substrings,
because a trigger that peeks is worse than a bad trigger: it looks like a good
one.

\textbf{What the detector cannot see.} Measured over 120 seeds per family at the
frozen calibration under the no-alert condition, the firing rate is $0.99$ for a
demand increase, $1.00$ for a demand decrease, $0.69$ for the temporary pulse,
$0.03$ for the lead-time shift, $0.03$ for the lost-shipment burst and $0.77$ for
the compound event. The two supply-only families sit at the null rate \emph{by
construction}, since a demand-side detector cannot see a lead-time shift or a
lost shipment. This is not a tuning failure to be fixed: it is what makes the
alert channel the only early signal for our running example.

\textbf{The alert bank.} The text channel is 180 short synthetic operational
templates, 60 per split and 10 per family within each split, of which four per
family-split cell are accurate, two overstated, two ambiguous and two operational
distractors, with wording families held out across splits and controlled slots
instantiated per episode. Template identity stays on the analysis side, because
language-level uncertainty must treat the template as a sampling unit and not the
rendered string, and synthetic messages are preferred because timing and
information content are controllable: a message must not carry information that
could not have existed when it is delivered. Two independent raters score every
template on five criteria, and leakage is not negotiable by consensus, since one
below-threshold score from either rater removes it. \pending{inter-rater
agreement and templates pruned for leakage} Each message also carries a hidden
\alertspec{}, and feeding that straight to the compiler is the perfect-parsing
upper bound. Every trajectory is crossed with the identical exogenous realization
under four information conditions, and the four cells of one seed are
\emph{paired replays of one independent unit}.
